% ---------- 扉页 ----------
\begin{titlepage}
\centering
\vspace*{0.24\textheight}
{\zihao{-4}\color{black!60}现代汉语选编与重述\par}
\vspace{1.1em}
{\zihao{1}\heiti 刚体运动理论\par}
\vspace{1.5em}
{\zihao{4}莱昂哈德·欧拉 \quad 著\par}
\vspace{2.8em}
{\small\color{black!60}%
底本：\textit{Theoria motus corporum solidorum seu rigidorum}（E289，1765）\par
\vspace{0.3em}
英文参照：Ian Bruce 英译本\quad ·\quad 重述与批注：本系列编者\par}
\vfill
{\footnotesize\color{black!52}2026 年 11 月\quad ·\quad 据工作区 Markdown 文稿自动排版\par}
\end{titlepage}

% ---------- 版本与读法说明页 ----------
\thispagestyle{empty}
\vspace*{0.10\textheight}
\begin{fullwidth}
\small
{\heiti\zihao{-4} 关于本册\par}\vspace{0.8em}
本册由 31 篇 Markdown 文稿自动排版合成，含总览、引言六章、正文十九章、附录五章。

\vspace{0.5em}
{\heiti\zihao{-5} 读法\par}\vspace{0.4em}
\begin{itemize}\small
  \item 正文为欧拉思想路线的现代汉语重述；
  \item 边栏楷体【批注】为编者补充（直观解释、现代对照、延伸入口），随文就位；
  \item 章前【选编说明】交代本章相对底本的取舍；
  \item 图题置于图下居中；全书图号连续编排；
  \item 点击目录、正文中的章节互引可跳转。
\end{itemize}

\vspace{0.5em}
{\heiti\zihao{-5} 底本\par}\vspace{0.4em}
Euler, L., \textit{Theoria motus corporum solidorum seu rigidorum}, E289, 1765. 英文参照为 Ian Bruce 英译（17thcenturymaths.com）及随附 \texttt{english\_translation/} 文本；现代对照见 Landau--Lifshitz 与 Goldstein 等教材，详见总览文末书目。
\end{fullwidth}
\clearpage

% ---------- 目录 ----------
\renewcommand{\contentsname}{目\quad 录}
\tableofcontents
